%!TEX program = uplatex
\RequirePackage{plautopatch}
\documentclass[uplatex,dvipdfmx]{jlreq}
\usepackage{amsmath,amssymb}

\pagestyle{empty}

\begin{document}

%================================================
% 講演1
%================================================

\begin{center}
{\large\bfseries 粘性解への招待}
\end{center}

\begin{flushright}
\textbf{石井 仁司（都立大・理）}
\end{flushright}

粘性解は偏微分方程式の弱解の一概念である。１９８１年に Crandall と Lions によって導入された。当初のもくろみは Hamilton-Jacobi 方程式に対する自然な大域解を如何に定義するかということであった。そこでは一階の微分方程式である Hamilton-Jacobi 方程式に対して通常の意味では微分可能でない解を考える必要があり、しかもその解は存在と一意性という点で十分によい性質を持つという要求があった。粘性解はこの要請に答えるものとして導入された。それはまた Kruzkov と Douglis による Hamilton-Jacobi 方程式のそれまでの研究を発展させるというものであり、さらに Kruzkov による保存則 (conservation laws) に対する弱解である entropy 解の Hamilton-Jacobi 方程式に対応する弱解は何かという問に答えるものであった。時代背景としては、非線形半群理論の具体的な偏微分方程式への応用を追求するという機運の下でなされた。

粘性解の導入とほとんど同時に、その一つの本質は最大値原理にあるという指摘により、その定義は理解し易いものに整理された。従って、対象となる方程式の自然なクラスは単独二階（退化）楕円型方程式であることが認識される。また、Schwartz の超関数と比すとき、超関数では「部分積分による微分の試験関数への移行」という方法を用い、粘性解では「最大値原理による微分の試験関数への移行」という方法を用いて、通常の意味で微分不可能な関数に微分概念を導入すると言える。

これまでの研究により、粘性解の重要な基本性質として、解の一意性あるいは解の比較が、Dirichlet 問題、Neumann 問題等に対して広いクラスの単独退化楕円型方程式に対して解の滑らかさの仮定なしに成立することや解が広義一様収束の位相における極限操作に関して安定性であることが知られている。もちろん、安定性は解の存在を示す上で重要である。連続関数の空間でコンパクトな近似解の列を何らかの方法で構成すれば、部分列の極限として粘性解が得られる。これまでに得られている比較定理が粘性解理論の中核を成している由だが、解の一意性（比較）の問題は永遠の課題に思われる。現在までに得られている比較定理に関する簡単なサーベイを通して、粘性解の理論の解説を行いたいと考えている。

粘性解のこれまでに知られている応用は概ねつぎの４種類に分けられる。

(1) 最適制御と微分ゲーム  
(2) 楕円型方程式あるいは放物型方程式に対する特異摂動  
(3) 超曲面の時間発展へのレベルセット・アプローチ  
(4) $L^\infty$ における変分法  

（１）については、Hamilton-Jacobi 方程式は最適制御と微分ゲームに対する値関数の満たすべき方程式として現れる。この値関数は通常せいぜい Lipschitz 連続程度の滑らかさしか持たない。

（２）については、広義一様収束に関する粘性解の安定性により、特異摂動問題に対して多くの応用を持つ。

（３）については、１９８０年代の終わりに粘性解の新しい応用として Chen-Giga-Goto と Evans-Spruck により最初に研究された。儀我美一氏の精力的な貢献により、理論と応用の両面で現在最も進展目ざましい。

（４）については、１９６０年代の Aronsson の研究に遡ることができるが、粘性解の導入により初めて解析の土台が与えられたといえる。

これらの応用の幾つかを簡単に紹介したいと考えている。

\bigskip

参考文献

[1] M. G. Crandall, H. Ishii, P.-L. Lions,  
User's guide to viscosity solutions of second order partial differential equations,  
Bull. AMS 27 (1992), 1--67.

[2] 石井仁司，粘性解の応用，数学 47 (1995), 97--110.

\newpage

%================================================
% 講演2
%================================================

\begin{center}
{\large\bfseries 粘性解という弱解の幾何学的側面}
\end{center}

\begin{flushright}
\textbf{儀我 美一（北大・理）}
\end{flushright}

必ずしも通常の意味で微分できない関数に対して解の概念を拡張することによって偏微分方程式の解とみなす考え方には１００年以上の長い歴史があります。

例えば電磁気学でおなじみのポアソン方程式やラプラス方程式の境界値問題の解を変分原理によって構成することを考えましょう。雑にいえばディリクレ積分という（汎）関数を最少にするような関数をみつける問題を解くことになります。

汎関数に下限があるとしましょう。直接法という方法では、ディリクレ積分の下限に近づいていく関数の列 --- 最小化列という --- をつくり（この列はいつも存在します）その極限として最小解を構成します。

ここで最小化列に何らかのコンパクト性があれば少なくとも部分列は収束します。もし汎関数がその収束に関して下半連続であれば、収束先こそ最小解です。

ここでコンパクト性をいうためにはできるだけ弱い位相で考えたいし、下半連続性をいうにはできるだけ強い位相で考えたくなるという問題があります。

うまく適当な位相がみつかってもその位相での収束極限は必ずしももとの方程式を記述するのに必要な階数だけ微分可能にはなりません。そこでどうしても解の概念を弱める必要が出てきました。これがいわゆる弱解の概念の始まりです。

弱解の概念を正しく捉えるには関数空間の概念が必要で、そのため関数解析が進展し、超関数論も生まれました。

しかし非線形方程式では事情が異なります。バーガーズ方程式やハミルトン・ヤコビ方程式のような非線形方程式ではデータが滑らかでも有限時間で解の滑らかさが失われます。

そこで１９８０年代前半に Crandall と Lions により導入された弱解の概念 --- 粘性解 --- が重要になります。

この理論は石井仁司先生により整備され、現在では大変理解しやすい理論となっています。

その後、平均曲率流方程式の弱解の構成に粘性解の理論が用いられると（Chen-Giga-Goto, Evans-Spruck）幾何解析の立場からも注目されるようになりました。

ここでは粘性解を幾何学的側面から見直しながら、曲面の方程式の弱解の構成法である等高面の方法を振り返ってみたいと思います。

\bigskip

参考文献

[1] 儀我美一・陳蘊剛，動く曲面を追いかけて，日本評論社 (1996)

[2] 儀我美一，曲面の発展方程式における等高面の方法，数学 47 (1995), 321--340

\newpage

%================================================
% 講演3
%================================================

\begin{center}
{\large\bfseries 「半連続粘性解」について}
\end{center}

\begin{flushright}
\textbf{小池 茂昭（埼玉大・理）}
\end{flushright}

儀我・長井氏によって粘性解理論の多彩な応用が本集会で紹介されると思われますが、私の講演ではもう一度原点に戻ってコントロールへの応用を念頭において議論の再点検を行いたいと思います。

石井氏の概説でも述べられる通り、粘性解の理論の長所として『期待される』解の一意性を保証している点が挙げられます。これは比較定理によって導かれます。

しかしこの比較定理は同時に解の連続性まで導いてしまいます。つまり期待される解が不連続の場合に粘性解の理論は適用できないことになります。

この困難を解決したのが Barron と Jensen による研究です。

この講演では Barron-Jensen の補題を解説するとともに、この概念の正当性を逆向き動的計画原理を用いて説明します。

\bigskip

参考文献

[1] E.N. Barron \& R. Jensen,  
Semicontinuous viscosity solutions for Hamilton-Jacobi equations with convex Hamiltonians,  
Comm. PDE 15 (1990), 1713--1742.

\newpage

%================================================
% 講演4
%================================================

\begin{center}
{\large\bfseries 粘性解と制御理論の関わり}
\end{center}

\begin{flushright}
\textbf{長井 英生（阪大・基礎工）}
\end{flushright}

粘性解の概念はその誕生当初から制御問題と深い関わりを持って生まれました。

P. L. Lions は 80年代初頭に来日した際、Krylov の著書  
``Controlled Diffusion Processes'' を自身の ``Bible'' と呼んだと言われています。

一様楕円性から退化楕円性へ、超関数解から粘性解へという飛躍は制御理論とも深く関係しています。

80年代後半には数学者と工学者の研究が一時的に離れる時期がありましたが、90年代に入り $H_\infty$ 制御理論の成熟とともに再び関係が深まりました。

この講演ではその歴史的背景を説明したいと思います。

\end{document}